% Minimal retained preamble for lane-p2 cards taken from
% sources/logic-fl02-conservation-of-ramsey-s-theorem-f/pinned-original-source.tex
% (Conservation of Ramsey's theorem for pairs, arXiv:2402.11616v3).
%
% Copied verbatim from the author's own file:
%   line 2      : \documentclass{article}
%   lines 16-25 : the author's theorem declarations
%   lines 46-88 : the author's macro definitions for the reverse-mathematics principles
%   line 91     : \card
%   lines 104-107: the natural sums and products, which rely on mathabx (source line 103)
%   line 116    : \qt
%
% Omitted author preamble lines and why (nothing was renamed or redefined):
%   line 3 (graphicx), 7 (url), 9 (cancel): no figure, url or cancel use in the retained spans;
%   line 6 (cleveref): the assembler adds its own cleveref machinery, a second load clashes;
%   line 8 (\usepackage[utf8x]{inputenc}): inputenc with the utf8x option fights the forced
%       xelatex engine;
%   lines 28-37 (the rep@theorem machinery and \newreptheorem{maintheorem}): the retained spans
%       use no rep-theorem environment;
%   lines 40-43 (xcolor and the three author comment macros): not used by the retained spans;
%   lines 93-99 (the LS2 font encoding block for \smashtimes based on stix): the host TeX Live has
%       no stix directory and the author himself never uses the result -- line 105 maps \ntimes to
%       \dottimes, and line 110 shows \smashtimes already commented out in the source.
% mathabx is retained because \dotplus, \dottimes and \dotdiv come from it.
\documentclass{article}
\usepackage{amsmath,amssymb,amsthm}
\usepackage{mathabx}
\usepackage{cancel}
\usepackage{cleveref}
\usepackage{url}

\newtheorem{theorem}{Theorem}
\numberwithin{theorem}{section}
\newtheorem{maintheorem}[theorem]{Main Theorem}
\newtheorem{lemma}[theorem]{Lemma}
\newtheorem{question}[theorem]{Question}
\newtheorem{proposition}[theorem]{Proposition}
\newtheorem{statement}[theorem]{Statement}
\newtheorem{remark}[theorem]{Remark}
\newtheorem{definition}[theorem]{Definition}
\newtheorem{corollary}[theorem]{Corollary}

\newcommand{\RCA}{\mathsf{RCA}}
\newcommand{\WKL}{\mathsf{WKL}}
\newcommand{\SRT}{\mathsf{SRT}}
\newcommand{\EM}{\mathsf{EM}}
\newcommand{\SEM}{\mathsf{SEM}}
\newcommand{\ADS}{\mathsf{ADS}}
\newcommand{\SCAC}{\mathsf{SCAC}}
\newcommand{\SADS}{\mathsf{SADS}}
\newcommand{\COH}{\mathsf{COH}}
\newcommand{\ACA}{\mathsf{ACA}}
\newcommand{\qvdash}{\operatorname{{?}{\vdash}}}
\newcommand{\nqvdash}{\operatorname{{?}{\nvdash}}}
\newcommand{\BSig}{\mathsf{B}\Sigma^0}
\newcommand{\ISig}{\mathsf{I}\Sigma^0}
\newcommand{\LSig}{\mathsf{L}\Sigma^0}
\newcommand{\BPi}{\mathsf{B}\Pi^0}
\newcommand{\IPi}{\mathsf{I}\Pi^0}
\newcommand{\LPi}{\mathsf{L}\Pi^0}
\newcommand{\IDel}{\mathsf{I}\Delta^0}
\newcommand{\BCDe}{\mathsf{BC}\Delta^0}
\newcommand{\WF}{\mathsf{WF}}
\newcommand{\RT}{\mathsf{RT}}
\newcommand{\BME}{\mathsf{BME}}
\newcommand{\TT}{\mathsf{TT}}
\renewcommand{\O}{\mathcal{O}}
\renewcommand{\L}{\mathcal{L}}
\newcommand{\C}{\mathcal{C}}
\newcommand{\R}{\mathcal{R}}
\renewcommand{\S}{\mathcal{S}}
\newcommand{\T}{\mathcal{T}}
\renewcommand{\P}{\mathcal{P}}
\newcommand{\M}{\mathcal{M}}
\newcommand{\Nc}{\mathcal{N}}
\renewcommand{\P}{\mathcal{P}}
\newcommand{\Q}{\mathcal{Q}}
\newcommand{\NN}{\mathbb{N}}
\newcommand{\PP}{\mathbb{P}}
\newcommand{\uh}[0]{{\upharpoonright}}
\newcommand{\card}{\operatorname{card}}
\newcommand{\nplus}{\dotplus}
\newcommand{\ntimes}{\dottimes}
\newcommand{\nminus}{\dotdiv}
\newcommand{\nsum}{\dot{\sum}}
\def\qt#1{``#1''}
\title{Ramsey's theorem for pairs preserves a finite level scheme without controlling the jump}
\author{Quentin Le Houérou, Ludovic Levy Patey, Keita Yokoyama}
\date{arxiv:2402.11616v3; CC BY 4.0}
\begin{document}
\maketitle
\noindent\textit{Editorial note.} The counted unit is the article's own Corollary cor:rt22-sig2 at source lines 1540-1547 together with the proof environment that begins on the next line (1548-1553); the pairing is the ordinary adjacent one.\par\medskip
\noindent\textit{Editorial note.} The retained context is the statement of Proposition prop:rt22-sig2 at source lines 1506-1514, which the proof applies; that proposition is counted by a sibling card of this lane and the counted intervals are disjoint.\par\medskip
\noindent\textit{Editorial note.} The counted proof carries the author citation of Belanger. The entry is supplied as a task-directory fragment holding it exactly as the bibtex run had formatted it, because the bibtex backend cannot be used here: the harness runs both xelatex passes before bibtex, so a generated bbl is never read back and the citation would stay undefined in the final pass.\par\medskip
\noindent\textit{Editorial note.} The author has left a commented-out alternative first step at source lines 1550-1551; it is commented out and is not author text for our purposes, and it is not retained.\par\medskip
\noindent\textit{Editorial note.} The retained preamble reproduces verbatim the author's document class at source line 2, his theorem declarations at source lines 16-25, his macro definitions at source lines 46-88 and 91, and his natural-operation macros at source lines 104-107, and it retains mathabx and cleveref; the stix-dependent LS2 block and the inputenc load with the utf8x option are not retained.\par\medskip
\noindent\textit{Editorial note.} The retained preamble also loads the url package because the attached bibliography entry carries a url.\par\medskip
\noindent\textit{Allowed standard prerequisites (editorial).}\begin{itemize}
\item The finite-level proposition of this section. For a countable topped model of RCA\_0 plus the scheme at level n minus one, every colouring and every auxiliary set above the jump of Y, there is a homogeneous set preserving the scheme at level n and controlling the jump; retained with this card at source lines 1506-1514.. The counted corollary is exactly this proposition with the auxiliary set supplied internally, so the statement must be available to see what is dropped.
\item Omega-extension for finite levels of induction and collection. Every countable topped model of RCA\_0-star plus the scheme at level n minus one can be omega-extended into a model of the same theory plus weak Konig's lemma, a result attributed to Harrington, Simpson-Smith and Hajek and recorded by Belanger; the citation is carried in the proof and its entry is attached with this card.. It replaces the external auxiliary set of the proposition by a set P obtained inside the model, which is the only new ingredient of the corollary.
\item RCA\_0-star and the jump of the topped model. RCA\_0-star is the fragment of RCA\_0 in which the extension step above is stated, and the model extended by its own jump satisfies that fragment at level n minus one because the scheme is assumed for the original model.. The proof applies the extension result to the model extended by its jump, which is why the corollary needs no external auxiliary set.
\end{itemize}\medskip
\begin{proposition}\label[proposition]{prop:rt22-sig2}
Fix~$n \geq 2$ and let $T_n$ be either $\ISig_n$, or $\BSig_{n+1}$.
Consider a countable model $\M = (M,S) \models \RCA_0 + T_{n-1}$ topped by a set~$Y \in S$. Then, for every coloring~$f : [M]^2 \to 2$ in~$S$ and every set~$P \gg Y'$ such that $\M[P] \models \RCA_0$, there exists $G \subseteq M$ such that 
\begin{enumerate}
    \item $G$ is an $M$-infinite $f$-homogeneous set ;
    \item $P \geq_T (G \oplus Y)'$ ;
    \item $\M[G] \models \RCA_0 + T_n$.
\end{enumerate}
\end{proposition}

\begin{corollary}\label[corollary]{cor:rt22-sig2}
Fix~$n \geq 2$ and let $T_n$ be either $\ISig_n$, or $\BSig_{n+1}$.
Consider a countable model $\M = (M,S) \models \RCA_0 + T_n$ topped by a set~$Y \in S$. Then, for every coloring~$f : [M]^2 \to 2$ in~$S$, there exists $G \subseteq M$ such that 
\begin{enumerate}
    \item $G$ is an $M$-infinite $f$-homogeneous set ;
    \item $\M[G] \models \RCA_0 + T_n$.
\end{enumerate}
\end{corollary}

\begin{proof}
By Harrington; Simpson–Smith; Hajek (see Belanger~\cite{Blanger2022ConservationTF}), every countable topped model of $\RCA_0^* + T_{n-1}$ can be $\omega$-extended into a model of~$\RCA_0^* + T_{n-1} + \WKL$. Since~$\M[Y'] \models \RCA_0^* + T_{n-1}$, there is some $P \subseteq M$ such that $P \gg Y'$ and $\M[P] \models \RCA_0^* + T_{n-1}$.
% By Harrington (see~\cite[Theorem IX.2.1]{simpson_2009}), every countable model of~$\RCA_0$ can be $\omega$-extended into a model of~$\WKL_0$.
% Therefore, there exists a set~$P \subseteq M$ such that $P \gg Y'$ and $\M[P] \models \RCA_0$.
The conclusion follows from~\Cref{prop:rt22-sig2}.    
\end{proof}

\begin{thebibliography}{99}
\bibitem{Blanger2022ConservationTF}
David~R. B{\'e}langer.
\newblock Conservation theorems for the cohesiveness principle.
\newblock 2022.
\newblock Available at \url{https://arxiv.org/abs/2212.13011}.
\end{thebibliography}
\end{document}
